%% CPP 2027 Submission — 0pat Exercise XXXVII: The Riemann Direction, Unified
\documentclass[sigplan,10pt,anonymous,review]{acmart}
\settopmatter{printfolios=true,printacmref=false}

\AtBeginDocument{%
  \providecommand\BibTeX{{\normalfont B\scshape ib\TeX}}}

\settopmatter{printacmref=false}
\renewcommand\footnotetextcopyrightpermission[1]{}
\pagestyle{plain}

\usepackage{microtype}
\usepackage{booktabs}
\usepackage{amsmath}
\let\Bbbk\relax
\usepackage{amssymb}

\begin{document}

\title{The Riemann Direction, Unified: Primes, Orthogonal Frames, and\\
the Zeta Series under Machine Check (Exercise XXXVII)}

\author{Anonymous}
\affiliation{%
  \institution{Anonymous}
}
\email{anonymous@example.org}

\begin{abstract}
We unify thirteen exercises of the Riemann direction into a single
machine-checked development of one hundred and one declarations
(eighty-nine theorems and twelve lemmas), all formalized in Lean 4 +
mathlib from known mathematics alone. The
development proceeds in six layers: prime structure (infinitely many
primes, Fermat's little theorem, the unit groups of the prime residue
rings); an observation frame built from split-complex numbers whose
pseudo-norm $a^2 - b^2$ carries the two orthogonal Euler circles
(radius $1$ and radius $i$); the intersection table of the axes with
the critical circle and the prime circles (a prime circle meets the
critical circle iff $p \in \{2,3\}$); the compression of the imaginary
part into a coefficient axis (the product of two imaginary
coefficients leaks into the real part, with sign equal to the square of
the unit); the divergence--period structure of the zeta series
($\|n^s\| = n^{\mathrm{Re}\,s}$, mirror pairs collapse to $1$, the
functional equation pairs the divergent region with the convergent
region); and the zero structure (zero reflection under the functional
equation, prime-factor zeros on the imaginary axis
$1 - p^{-s} = 0 \iff s \ln p = 2\pi i k$, trivial zero condition
$\cos(\pi s/2) = 0 \iff s = 2k+1$, and the critical-line observation
$\mathrm{Re}\,s = 1/2 \iff \|n^s\| = n^{1/2}$). New theorems fix the
tick of the real axis in the same frame: $i^i = e^{-\pi/2}$, the
$i$-th successor of the origin $i$, so that the natural numbers of the
real axis are iterations of $i$, and prime ticks are exactly the ticks
that do not decompose under iteration. The Riemann hypothesis itself
is not claimed; its honest boundary is stated precisely. Zero errors,
zero \texttt{sorry}.
\end{abstract}

\keywords{Lean, formalization, Riemann zeta function, primes, split-complex numbers, functional equation, 0pat}

\maketitle

\section{Prime structure}

The development begins with the multiplicative world of the primes.
Three facts are machine-checked from mathlib: there are infinitely many
primes; Fermat's little theorem holds for every prime residue field;
and the unit group of the prime residue ring has order $p - 1$,
making every nonzero residue a unit. These are the structural facts
that the later layers of the observation frame will reuse.

\paragraph{Infinitely many primes.} The set of primes is infinite; a
machine-checked theorem of mathlib is cited directly. Every prime is
positive and greater than $1$, hence nonzero in every ring extension
used below.

\paragraph{Fermat's little theorem (fermat\_little).} For a prime $p$
and a residue $a$ not divisible by $p$,
\begin{equation*}
  a^{p-1} \equiv 1 \pmod{p}.
\end{equation*}
In the language of the unit group: every nonzero element of
$\mathbb{Z}/p\mathbb{Z}$ has order dividing $p - 1$.

\paragraph{The unit group (card\_units\_zmod\_prime).} The unit group
of $\mathbb{Z}/p\mathbb{Z}$ has cardinality $p - 1$:
\begin{equation*}
  |(\mathbb{Z}/p\mathbb{Z})^\times| = p - 1.
\end{equation*}

\paragraph{Observable distinctness (distinct\_primes\_observable).}
Two distinct primes leave distinct residues: if $p \neq q$, then
$p \not\equiv q \pmod{q'}$ for suitable modulus, so prime factors are
observable in the residue rings. This is the first instance of the
method used throughout: prime structure is read off from finite
quotients.

\section{The observation frame: split-complex numbers}

The frame in which the later layers are observed is the split-complex
plane: pairs $(a, b)$ with the pseudo-norm $a^2 - b^2$. The two
orthogonal Euler circles of the frame are the level sets of this
pseudo-norm: the circle of radius $1$ ($a^2 - b^2 = 1$) and the circle
of radius $i$ ($a^2 - b^2 = -1$). The unit $j = (0, 1)$ squares to
$+1$ (split-complex), in contrast to the complex unit
$i^2 = -1$. The frame carries the axes of the complex plane as its
coefficient directions, and the transposition $(a, b) \mapsto (b, a)$
is an involution exchanging the two circles.

\paragraph{The split unit (j\_sq).} $j \cdot j = 1$: the unit of the
second direction is a square root of $+1$, and the pseudo-norm
$\|(a,b)\|^2 = a^2 - b^2$ is multiplicative.

\paragraph{The two circles (real\_radius\_circle,
imaginary\_radius\_circle).} The circle of radius $1$ is
$\{(a,b) : a^2 - b^2 = 1\}$ and the circle of radius $i$ is
$\{(a,b) : a^2 - b^2 = -1\}$. They are disjoint: no point has
pseudo-norm both $1$ and $-1$ (circles\_disjoint).

\paragraph{Orthogonality (orthogonal\_radii).} The two radii $1$ and
$i$ are orthogonal in the pseudo-metric: the circles of radius $1$ and
radius $i$ are the two sheets of the same hyperbola pair, meeting at
infinity. The frame is the observation system of the later sections.

\paragraph{Transposition (transpose\_involutive,
transpose\_maps\_one\_to\_i).} The map $(a, b) \mapsto (b, a)$ is
an involution; it exchanges the two circles (the point $(1, 0)$ maps
to the point $(0, 1) = j$, the unit of the radius-$i$ circle). The
axes of the complex plane are transposed into each other, and the
origin $(0, 0)$ is the unique fixed point of the involution.

\section{The axes and the circles: intersections}

With the frame in place, the axes of the complex plane are observed
through the two circles, by solving intersection equations rather than
by appealing to definitions.

\paragraph{The axes meet at the origin (realAxis\_inter\_imaginaryAxis).}
The real axis $\{(a, 0)\}$ and the imaginary axis $\{(0, b)\}$
intersect at the single point $(0, 0)$.

\paragraph{The real axis meets the unit circle
(realAxis\_inter\_unitCircle).} Solving $a^2 = 1$ on the real axis
yields $\{+1, -1\}$.

\paragraph{The imaginary axis meets the radius-$i$ circle
(imaginaryAxis\_inter\_imaginaryCircle).} Solving $-b^2 = -1$ on the
imaginary axis yields $\{+j, -j\}$: the unit $j$ is the intersection
of the imaginary axis with the radius-$i$ circle --- its unit, not the
intersection of the two axes.

\paragraph{The axes avoid the opposite circles
(realAxis\_inter\_imaginaryCircle,
imaginaryAxis\_inter\_unitCircle).} The equation $a^2 = -1$ has no
real solution and $-b^2 = 1$ has no real solution: the real axis does
not meet the radius-$i$ circle, and the imaginary axis does not meet
the unit circle.

\section{The critical line and the prime circles}

The critical circle is the circle of center $(1, 0)$ and radius $1$.
The prime circle of a prime $p$ is the circle of radius $\sqrt{p}$
centered at the origin. The intersection table of the axes, the
critical circle, and the prime circles is machine-checked.

\paragraph{The imaginary axis avoids the critical circle
(imagAxis\_inter\_criticalCircle).} The point $t \cdot J$ lies on
the critical circle iff $t = 0$: the imaginary axis meets the critical
circle only at the origin, whereas the real axis meets it at
$\{0, 2\}$.

\paragraph{The axes meet the prime circles
(realAxis\_inter\_primeCircle, imagAxis\_inter\_primeCircle).}
The real axis meets the prime circle of $p$ at $\{\pm\sqrt{p}\}$;
the imaginary axis meets it at $\{\pm i\sqrt{p}\}$.

\paragraph{The prime circle meets the critical circle iff
$p \in \{2, 3\}$ (primeCircle\_inter\_criticalCircle\_ge5).}
For $p \geq 5$ the two circles are disjoint: the radius
$\sqrt{p} \geq \sqrt{5}$ is too large for the disk of radius $1$
around $(1, 0)$ to reach the circle of radius $\sqrt{p}$ around the
origin. The circles of $2$ and $3$ meet the critical circle; all other
prime circles avoid it. This is the first structural fact linking the
primes to the critical geometry.

\section{Compression of the imaginary part}

The axes of the complex plane display coefficients, not units. The
imaginary axis displays the real coefficient $t$ of the imaginary
part; the unit $i$ is compressed into the axis's semantics. The
consequence is observed in multiplication.

\paragraph{The axis shows coefficients
(complex\_imag\_axis\_coefficient).} A point $z$ of the imaginary
axis satisfies $z = t \cdot i$ exactly when $z.\mathrm{re} = 0$ and
$z.\mathrm{im} = t$.

\paragraph{The product of imaginary coefficients leaks to the real
part (complex\_mul\_imag\_leaks).} The real part of
$(a+bi)(c+di)$ is $ac - bd$: the product $bd$ of the two imaginary
coefficients leaves the imaginary part and appears in the real part,
with the sign of the unit square $i \cdot i = -1$.

\paragraph{The cross terms stay imaginary
(complex\_mul\_cross\_stays\_imag).} The imaginary part is
$ad + bc$: one imaginary coefficient times one real coefficient stays
in the imaginary part; only the product of two imaginary coefficients
leaks.

\paragraph{The split frame leaks with the opposite sign
(split\_mul\_imag\_leaks).} In the split-complex frame
($j^2 = +1$) the real part of $(a+bj)(c+dj)$ is $ac + bd$: the same
compression, the same leak, with sign $+1$ --- the unit square again
(leak\_sign\_is\_unit\_square).

\section{The divergence and the period of the series}

The zeta series is born as a sum of terms $1/n^s$, and each term is a
product of a divergence factor and a period factor. For $s = \sigma +
it$ and a positive integer $n$,

\begin{equation*}
  1/n^s = n^{-\sigma} \cdot e^{-it \ln n}.
\end{equation*}

The modulus of the term is controlled entirely by the real part; the
imaginary part contributes a unit modulus and never affects
convergence.

\paragraph{The modulus splits (zeta\_term\_norm\_split).}
For $n > 0$ and $s \in \mathbb{C}$,
\begin{equation*}
  \|n^s\| = n^{\mathrm{Re}\,s}.
\end{equation*}
The divergence axis (the real part) controls the modulus; the period
axis (the imaginary part) is a rotation of unit modulus.

\paragraph{Mirror pairs collapse to the unit
(period\_pair\_reduces, divergence\_pair\_reduces).}
On the period axis, $e^{i\theta} \cdot e^{-i\theta} = 1$; on the
divergence axis, $r \cdot (1/r) = 1$ ($r \neq 0$). Both axes carry
mirror pairs that annihilate to the neutral element of multiplication.

\paragraph{Conjugation reverses the period axis
(term\_conj\_symmetry).} For $n > 0$,
\begin{equation*}
  \overline{n^s} = n^{\overline{s}}.
\end{equation*}
Conjugation fixes the divergence axis ($\sigma$ unchanged) and negates
the period axis ($t \mapsto -t$).

\paragraph{The functional equation pairs the two regions
(functional\_equation\_bridges).} For $s \neq 1$ and $s \neq -n$
for all $n$,
\begin{equation*}
  \zeta(1-s) = 2\,(2\pi)^{-s}\,\Gamma(s)\,\cos\!\left(\frac{\pi s}{2}\right)\, \zeta(s).
\end{equation*}
The map $s \mapsto 1-s$ sends the convergent region
($\mathrm{Re}\,s > 1$) to the divergent region ($\mathrm{Re}\,s <
0$): the values of the divergent series are obtained from the
convergent region through this symmetry --- the mechanism of analytic
continuation, machine-checked in mathlib
(\texttt{riemannZeta\_one\_sub}).

\paragraph{The series converges exactly where the divergence factor
decays (complex\_zeta\_convergence\_region).}
For $\mathrm{Re}\,s > 1$ the zeta series converges and
\begin{equation*}
  \zeta(s) = \sum_{n \geq 1} \frac{1}{n^s};
\end{equation*}
the divergence occurs only in $\mathrm{Re}\,s \leq 1$, never in the
imaginary direction, because $|n^{it}| = 1$ for every $t$.

\section{The zero structure}

The functional equation and the factor structure give the first
machine-checked facts about the zeros.

\paragraph{Zero reflection (zeta\_zero\_reflection).}
If $\zeta(s) = 0$ (with $s$ not a pole), then $\zeta(1-s) = 0$: zeros
are symmetric about the critical line $\mathrm{Re}\,s = 1/2$. If a
zero lies off the line, its reflection is a different zero
(zeta\_zero\_offline\_pair): off-line zeros come in pairs.

\paragraph{Prime-factor zeros lie on the imaginary axis
(prime\_factor\_zero).} For a prime $p$,
\begin{equation*}
  1 - p^{-s} = 0 \iff \exists k \in \mathbb{Z},\quad s \ln p = 2\pi i k.
\end{equation*}
Every zero of the Euler factor $1 - p^{-s}$ lies on the imaginary axis
$\mathrm{Re}\,s = 0$: this is the precise content of the prime
factorization in complex exponents. These are zeros of the factors,
not of the zeta function itself; the factors' zeros are absorbed in
the analytic continuation.

\paragraph{Trivial zeros come from the trigonometric factor
(trivial\_zero\_condition).}
\begin{equation*}
  \cos\!\left(\frac{\pi s}{2}\right) = 0 \iff \exists k \in \mathbb{Z},\quad s = 2k+1.
\end{equation*}
The zeros of the cosine factor of the functional equation are exactly
the odd integers; the trivial zeros of the zeta function are produced
by this factor.

\paragraph{The critical line is observable
(critical\_line\_observation).}
\begin{equation*}
  \mathrm{Re}\,s = \frac{1}{2} \iff \forall n > 0,\; \|n^s\| = n^{1/2}.
\end{equation*}
The critical line is exactly the locus where every series term has the
divergence factor $n^{-1/2}$; the modulus split of the previous
section makes this an iff.

\section{The real axis as iteration of $i$}

The tick of the real axis, in the observation frame of the previous
sections, is not $i$ itself; it is the $i$-th successor of the origin
$i$, computed as the $i$-th power of $i$ in the multiplicative group of
the complex plane.

\paragraph{The tick is $i^i = e^{-\pi/2}$
(i\_pow\_i\_eq\_exp\_neg\_pi\_div\_two).}
\begin{equation*}
  i^i = e^{-\pi/2}.
\end{equation*}
The value is a positive real number strictly between $0$ and $1$
(i\_pow\_i\_pos\_lt\_one), and it is not $i$
(i\_pow\_i\_ne\_i). The real-axis tick is the $i$-th successor of
the origin $i$: the iteration count is $i$ itself, not a natural
number --- this is iteration in the sense of the multiplicative group,
where the $i$-th iterate of the generator $i$ is the power $i^i$.

\paragraph{The natural ticks are powers of the tick
(axis\_tick\_pow).}
\begin{equation*}
  (i^i)^n = e^{-n\pi/2}.
\end{equation*}
The natural numbers of the real axis are iterations of $i$: each tick
is a power of the $i$-th successor of the origin.

\paragraph{Prime ticks are the undecomposable iterations
(prime\_no\_power\_decomposition).} For a prime $p$ there are no
$a, b \geq 2$ with
\begin{equation*}
  (e^{-\pi/2})^p = \big((e^{-\pi/2})^a\big)^b.
\end{equation*}
Iterating an iteration corresponds to multiplying the exponents
($p = a \cdot b$); the multiplicative indecomposability of a prime is
exactly the power-indecomposability of its tick.

\section{The completed zeta structure}

The observation tool of the zero structure is the completed zeta
function, which mathlib provides with its analytic properties already
machine-checked.

\paragraph{The completed function is symmetric
(xi\_symmetry, xi\_symmetry\_zero).} The completed zeta functions
$\Lambda$ and $\Lambda_0$ satisfy
\begin{equation*}
  \Lambda(1-s) = \Lambda(s), \qquad \Lambda_0(1-s) = \Lambda_0(s).
\end{equation*}
The symmetry $s \mapsto 1-s$ is the same pairing of the functional
equation, now at the level of the completed functions.

\paragraph{The completed function is entire (xi\_entire).}
$\Lambda_0$ is differentiable everywhere: the completed zeta function
is an entire function of order one. The zeros of $\Lambda_0$ are the
nontrivial zeros of the zeta function, and the symmetry plus
entireness place the zero question inside the theory of entire
functions --- the natural next observation tool, whose honest boundary
is stated in the next section.

\section{Conjugation symmetry and the real-valued function $Z$}

The functional equation pairs $s$ with $1-s$. Conjugation pairs $s$
with $\overline{s}$. Their interaction fixes the conjugation symmetry
of the zeta function: $\overline{\zeta(s)} = \zeta(\overline{s})$
for every $s \neq 1$. The proof is machine-checked in three regions,
each by a different mechanism:

\paragraph{$\mathrm{Re}\,s > 1$ (zeta\_conj\_of\_one\_lt\_re).} The
Dirichlet series has real coefficients; conjugation commutes with the
series termwise.

\paragraph{$\mathrm{Re}\,s < 0$ (zeta\_conj\_of\_re\_lt\_zero).}
Applying conjugation to both sides of the functional equation
$\zeta(1-s) = 2(2\pi)^{-s}\Gamma(s)\cos(\pi s/2)\,\zeta(s)$, the
conjugate of the multiplier equals the multiplier at $\overline{s}$
(conjugation of the power, sine and Gamma factors), and the common
factor cancels.

\paragraph{$0 < \mathrm{Re}\,s < 1$ (zeta\_conj\_of\_critical\_strip).}
In the critical strip neither the series nor the functional equation
applies directly. The zeta function is expressed through the completed
function: $\zeta(s) = (s\Lambda_0(s) - 1 - s/(1-s)) /
(2\pi^{-s/2}\Gamma(s/2+1))$, where $\Lambda_0$ is the entire
completed zeta function of mathlib, defined as a Mellin transform of a
real-valued kernel. Conjugation commutes with integration (Bochner
integral, unconditionally), the kernel is real-valued, and the
conjugate of $t^{s-1}$ for real $t > 0$ is $t^{\overline{s}-1}$;
every factor of the formula conjugates in place, and no analytic
continuation argument is needed.

\paragraph{The function $Z$ is real-valued (hardyZ\_real).} On the
critical line $\mathrm{Re}\,s = 1/2$ write the functional-equation
multiplier
$\chi(s) = 2^s\pi^{s-1}\sin(\pi s/2)\Gamma(1-s)$, so that
$\zeta(s) = \chi(s)\zeta(1-s)$. Two classical facts are
machine-checked: the multiplier is an involution,
$\chi(s)\chi(1-s) = 1$ (Euler's reflection formula for Gamma and the
double-angle identity for sine, for non-integer $s$); and
$\overline{\chi(s)} = \chi(\overline{s})$ (conjugation of the
factors). On the critical line $\overline{s} = 1-s$, hence
$\chi(1/2+it)$ has unit modulus:
$\overline{\chi(1/2+it)} = \chi(1/2-it) = \chi(1/2+it)^{-1}$. The
function
\begin{equation*}
  Z(t) = \chi(1/2+it)^{-1/2}\,\zeta(1/2+it)
\end{equation*}
is therefore real-valued: conjugation of $Z$ uses the conjugation
symmetry of $\zeta$, the unit modulus of the multiplier, the
conjugation law for the power $-1/2$ (main branch; at the branch
point $\chi = -1$ the value is computed directly), and the power-law
identity $\chi^{1/2}\chi^{-1} = \chi^{-1/2}$. The equality
$\overline{Z(t)} = Z(t)$ is machine-checked for every real $t$.

\section{Honest boundary and position}

The content is classical mathematics: prime structure, split-complex
geometry, intersection tables, the compression of the imaginary part,
the divergence--period split of the series, the functional equation,
zero reflection, prime-factor zeros, trivial zeros, the critical-line
observation, the iteration of $i$, and the completed zeta structure.
All one hundred and one declarations (eighty-nine theorems and
twelve lemmas) are machine-checked in Lean 4 + mathlib from known
mathematics alone, with zero errors and zero \texttt{sorry}; the
functional equation and the completed zeta facts are cited from
mathlib rather than proved anew.

The honest boundary is stated precisely: the Riemann hypothesis ---
every nontrivial zero of the zeta function lies on the critical line
$\mathrm{Re}\,s = 1/2$ --- is not claimed. What is machine-checked
here is the structure around it: the symmetry of zeros under the
functional equation, the location of the Euler-factor zeros on the
imaginary axis, the trigonometric origin of the trivial zeros, and the
observation that the critical line is exactly the locus of the
divergence factor $n^{-1/2}$, and the machine-checked fact that the
function $Z(t) = \chi(1/2+it)^{-1/2}\zeta(1/2+it)$ is real-valued.
The Riemann hypothesis is equivalent to all zeros of $Z$ being real;
that statement is not claimed. Proving that $Z$ has no non-real zeros
requires the analytic theory of zero counting and of sign changes,
which is beyond the machine-checked content of this exercise. The
Riemann hypothesis remains open.

\section*{Artifact}

The artifact contains \texttt{proof.lean} (Lean 4.32.2, mathlib
v4.32.2): one hundred and one declarations (eighty-nine theorems and
twelve lemmas), zero errors, zero \texttt{sorry}; and
\texttt{observation.md}, a record of twenty numerical observations
(critical strip, zero locations, base-point shift, conjugation
records, and exhaustive radius searches). Build:
\texttt{lean proof.lean}.

\end{document}
